% /root/monodromy, 2026-11-19. Complete independent parent proof and code audit passed.
\subsection{A symbolic mixed contact for an arbitrary translation}
\label{subsec:ss-uniform-contact}
\status{Proved local calculation with an exact universal polynomial
certificate; the parent independently read the complete proof and program.
The following subsection proves source smoothness. Other collisions
for odd kernels with injective coordinate projections are excluded in
Proposition~\ref{prop:ss-uniform-graph-mixed-contact}. Uniform
tangent-triple bridges are proved in
Proposition~\ref{prop:ss-uniform-tangent-triple}, and every odd-degree
full image is determined in Theorem~\ref{thm:ss-all-odd-full-groups}.}

Let $E:y^2+y=x^3$ over an algebraically closed field of characteristic
two, let $X=E^2$, and let $L=\mathcal O(3O)\boxtimes\mathcal O(3O)$.
Take a translation
\[
 g=((a,b),(c,d)),\qquad b^2+b=a^3,\quad d^2+d=c^3,
\]
whose two coordinates differ from $O$, and suppose first that
$(c,d)\notin E[3]$. Put $P=(O,O)$ and $Q=g$.

\begin{lemma}[A uniform local mixed-contact slice]
\label{lem:ss-uniform-local-contact}
There is a nonempty irreducible three-dimensional parameter family
of coefficient matrices $S$, on a dense open of which the following
hold. The section $s$ vanishes at $P,Q$ and is smooth at both points;
the two translated tangent directions coincide, with contact order
exactly two. One has $\det S\ne0$ and $Q_H(S)=0$, where
$Q_H(S)=\sum S_{ij}^2S_{2-i,2-j}$ is the source Hasse cubic.
There is an explicitly given section $u$ with
\[
                         M(s)u=s^{[2]},
\]
and the local residue differential represented by $u/s$ has unequal
values on the two translated tangent vectors. The assertions concern
the indicated two points and the global coefficient operator only;
the source is not asserted to be smooth elsewhere.
\end{lemma}
\begin{proof}
Use three independent parameters $F,T,U$, all nonzero, and adjoin $v$
satisfying
\begin{equation}\label{eq:ss-uniform-contact-AS}
 v^2+v+U^3+U(c^4+c)+F^2(T+a^4+a)+FT^2=0.
\end{equation}
This Artin--Schreier equation is irreducible over $k(F,T,U)$:
its right side has a pole of odd order three at $U=\infty$,
whereas a nonzero pole of $h^2+h$ has even order. Hence it defines
an irreducible three-dimensional family. Set
\[
\begin{aligned}
 C&=T+a,& E_0&=U+c^2,& D_0&=v+Uc+d^2+Fa^2,\\
 B&=Tc^2+E_0a,& A&=Td^2+D_0a+Fb,\\
 S&=\begin{pmatrix}A&B&C\\D_0&E_0&1\\F&0&0\end{pmatrix}.
\end{aligned}
\]
Here the section is expressed in the tensor basis $(1,x,y)$ on
both factors. Direct substitution gives
\[
 s(P)=s(Q)=0,\qquad ds(P)=(1,0),\qquad ds(Q)=(v,0),
 \qquad\det S=FTU.
\]
The first derivative at $P$ uses the $Y=1$ chart at $O$ and the
invariant local tangent coordinate. At $Q$ it uses the invariant
derivations $\delta x=1,\delta y=x^2$.
Moreover
\[
 Q_H(S)=D_0^2+D_0+E_0^3+C^2F+F^2C,
\]
which is precisely the left side of
\eqref{eq:ss-uniform-contact-AS} after the elliptic equations
are used. The parameter $v$ is generically nonzero, so both indicated
points are smooth.

Let $J$ reverse the three coordinates, put $R=JS^{[2]}J$, and define
\[
                       u=\frac{S R^t S}{FTU}.
\]
We regard the right side as the matrix of its nine coefficients.
The root-removed Cartier matrix, after multiplication on the output
by $J\otimes J$, is the Hessian of $\det S$.
The inverse-Hessian identity gives
$M(s)u=s^{[2]}$ when $Q_H(S)=0$; equivalently, for a universal matrix
$S$, every coefficient of
\[
 M(s)\operatorname{vec}(S R^t S)
                   -(\det S)s^{[2]}
\]
is divisible by $Q_H(S)$. This identity is also verified coefficient
by coefficient by the exact program below.
The values of the residue form on the common translated tangent
are $u(P)$ and $u(Q)/v$. They satisfy
\[
 u(P)=\frac{F}{FTU},\qquad
 \frac{u(Q)}{v}=\frac{T\bigl(v^2+F^2(a^4+a)\bigr)}{FTU}.
\]
Their difference is therefore
\begin{equation}\label{eq:ss-uniform-contact-value}
       \frac{F+T\bigl(v^2+F^2(a^4+a)\bigr)}{FTU}.
\end{equation}
It is generically nonzero. Indeed if its numerator vanished, the
Artin--Schreier equation would give
\[
 v=F/T+U^3+U(c^4+c)+F^2T+FT^2.
\]
Squaring this and substituting back into
$v^2+F^2(a^4+a)=F/T$ produces a polynomial in $U$ with nonzero
leading term $U^6$, which is impossible for the independent parameters.

We finally check the contact order using actual elliptic addition,
not merely tangent vectors. At $O$ write the $Y=1$ coordinates as
$(t,z)$ with $z+z^2=t^3$, hence $z=t^3+O(t^6)$. Since the section
has derivative $(1,0)$ at $P$ and its restriction to the second
elliptic fibre is $Fz$, its local branch through $P$ has first
coordinate of order at least three. Translation by $g$ therefore
fixes the first elliptic coordinate through order two. On the second
factor the group law gives
\[
 x((c,d)+(t,z))=c+t+c^2t^2+O(t^3),\qquad
 y((c,d)+(t,z))=d+c^2t+ct^2+O(t^3).
\]
On the first-coordinate fibre at $(a,b)$ the section equals
$T(d^2+c^2x+y)$. Its value along the translated source branch is
thus
\[
                         T(c^4+c)t^2+O(t^3).
\]
The coefficient is nonzero because $(c,d)\notin E[3]$: on this cubic
its finite flexes have $c\in\mathbf F_4$, and the corresponding
nine points including $O$ are exactly $E[3]$. This proves contact
order two and completes all stated local assertions.
\end{proof}

\begin{remark}[The other coordinates and the global consequences]
If the second coordinate is a nonzero three-torsion point but the
first is not, interchange the two factors. If both coordinates are
nonzero three-torsion points, product automorphisms preserving $O$
act transitively on these pairs, so the already checked contact03
witness supplies local mixed contacts after an automorphism.
These facts remove a local polynomial-identity obstruction for any
translation with both coordinates nonzero. Source smoothness is proved in
Lemma~\ref{lem:ss-uniform-source-smoothness}. Proposition~\ref{prop:ss-uniform-graph-mixed-contact} excludes all other
nonordinary identifications for every odd kernel with injective
coordinate projections. Proposition~\ref{prop:ss-uniform-tangent-triple}
supplies the remaining triple bridges, giving the odd-degree theorem
\ref{thm:ss-all-odd-full-groups}.
\end{remark}

\begin{remark}[Exact certificate]
The retained program
\path{agents/monodromy/code/supersingular_uniform_contact_identities.py}
works in the polynomial ring over $\mathbf F_2$, reducing only by the
two displayed elliptic equations. It checks all containment, gradient,
determinant and cubic identities; all nine universal inverse-Hessian
coefficients; both residue evaluations; the nonzero degree-six
obstruction to their equality; and the rational elliptic-addition
expansion through order two. Its output is
\path{agents/monodromy/outputs/supersingular_uniform_contact_identities.json}.
The completed run passed in $0.34270756500336574$ seconds. No geometric
existence assertion beyond the local lemma is inferred from this code.
\end{remark}
